% ******************************* Reply to examiner report ********************************

\chapter*{Replies to examiners' reports}

The examiners' reports, written in response to the original thesis submission, highlight many specific, necessary corrections to be made, many general suggestions for improvement and many comments on the overall structure and choice of subject matter. In order to deal with all of this feedback appropriately I have here extracted the individual suggestions from the examiners' report and the two preliminary reports and, over the following pages, I give individual responses to each piece of feedback. Many of the comments in the preliminary reports are duplicated in the ten comments of the general examiners' report; in theses cases I give a brief reply and direct attention back to the general report rather than copying-and-pasting entire paragraphs. All changes to the original manuscript are shown in {\Red{red}} in this submitted revision with the exception of the final chapter ``Conclusions'' (page~\pageref{chap:Conc}) since the entire chapter is a new inclusion and so is not printed in {\Red{red}} to make the reading easier. A `clean' manuscript, without any {\Red{red}} highlighting, and without this ``Replies'' section, is supplied separately. 

The feedback of the examiners has been very helpful in re-shaping this thesis into a stronger piece of work and the opportunity to do so has been greatly appreciated. Particular attention has been paid to properly stating the mathematical reasons behind using the power spectrum scaling exponent as an early warning signal indicator, even when the power spectrum of the time series in question does not exhibit true asymptotic power-law scaling (see comment~\ref{rep:gen:mathematical} of the general report, below): this new discussion now makes up most of chapter~\ref{chap:1D}. One criticism of both examiners was the choice of the tropical cyclone problem as an application of our EWS techniques to a geophysical system: this chapter remains in the thesis but care has been taken to address the concerns of the examiners in the replies below and in the manuscript. In particular, we carefully state the reasons behind this choice, address the obvious shortcomings, and, in the new conclusions chapter, we discuss other geophysical systems which may be studied in future work, and ways in which the study of the tropical cyclone problem might be improved if better data were available. 

We thank the examiners for the time taken to review this re-submitted manuscript and hope that the original concerns have now been addressed satisfactorily. 

\section*{Reply to the general examiners' report}
\label{rep:gen}

\begin{enumerate}
	\item \label{rep:gen:debate} \emph{Very little is reported of the debate on tipping points in the Earth system. This needs to be expanded in order to put the results into a broader context.}\\
	— A section in the introductory chapter now provides some examples of tipping points in the Earth system, followed by a discussion of the different techniques which have been used to detect or predict tipping points in some of these systems (see section~\ref{sec:intro:geophysical}, page~\pageref{sec:intro:geophysical}). %{\Red{Maybe more references in conclusions to put results `into context'.}}
	\item \label{rep:gen:threekinds} \emph{The three kinds of tipping behaviours need to be described more thoroughly.}\\
	— The three types of tipping, as set out by \cite{livina2011,ashwin2012}, are now explained more thoroughly in chapter~\ref{chap:intro} (see section~\ref{sec:intro:classification_of_TPs}, page~\pageref{sec:intro:classification_of_TPs}). Examples of each are also studied in chapter~\ref{chap:1D} (see section~\ref{subsec:1D:indicatorComparison:TypesOfTipping}, page~\pageref{subsec:1D:indicatorComparison:TypesOfTipping}) for signs or early warning signals, although we note that our EWS indicator techniques are not sensitive to `noise-induced' tipping.
	\item \label{rep:gen:longdatasets} \emph{A fundamental issue associated to using these [scaling] methods is the need for long datasets and large ensembles thereof. ... Instead, most signals investigated by Joshua are pretty short and the ensembles are rather small. The sensitivity of the results to the length of the signals and the size of the ensemble should be investigated.}\\
	— We now give a detailed study of the sensitivity of the PS indicator to time series length in section~\ref{subsec:1D:Exponents_as_EWS:PSsensitivity} (page~\pageref{subsec:1D:Exponents_as_EWS:PSsensitivity}) and conclude that the performance improves when longer time series are available. However, when attempting to find an early warning signal in a toy model we continue to use relatively short time series where the `window' in which the PS exponent is estimated may only be $\approx 100$ points at each time step. In these cases we wish to compare the PS indicator to the ACF1 indicator (which uses lag-1 autocorrelation) and so we are making the comparison particularly stark in order to highlight the differences in the two methods. The study in section~\ref{subsec:1D:Exponents_as_EWS:PSsensitivity} also shows that we can expect better performance of the PS indicator close to tipping points which are characterised by critical slowing down, which is exactly the situation is which we wish to use this technique anyway. We use the PS indicator not as a precise measure of scaling but as a sensor of the changes in the temporal organisation of the time series. In addition, the data we obtained for the study of tropical cyclones in chapter~\ref{chap:cyclones} is necessarily quite short and the ensembles necessarily quite small. These shortcomings are remarked upon in the new conclusion chapter (chapter~\ref{chap:Conc}, page~\pageref{chap:Conc}).
	Unfortunately, in many real systems gathered data may be either short/incomplete, or the dynamics of the process vary quickly in relation to the sampling rate of measurement; even in these circumstances our techniques may be useful, as shown in the case studies.
	The demonstration of the use of EWS indicator techniques with relatively short time series is, we consider, a prelude to chapter~\ref{chap:cyclones} in which this is tested. 
	\item \label{rep:gen:periodic} \emph{... Additionally the presence of periodic components can make these indicators (especially the DFA) useless. }\\
	— One hypothesis of this thesis is that the PS indicator is better-suited for EWS analysis in this scenario than other indicators, and in section~\ref{subsec:1D:Exponents_as_EWS:Periodicity} (page~\pageref{subsec:1D:Exponents_as_EWS:Periodicity}) we make a detailed study of this by comparing the abilities of the PS exponent, DFA exponent and lag-1 autocorrelation to accurately estimate the parameter of an AR(1) process from a time series when that process has been super-imposed with different periodic functions. We find that, indeed, the DFA exponent and autocorrelation are greatly affected by the introduced periodicity, so as to make their use in practical examples useless, whilst the PS exponent is almost completely unaffected. In the case of the sea level pressure data in our study of tropical cyclones (chapter~\ref{chap:cyclones}), we have attempted to remove the periodicity by two different methods (see section~\ref{subsec:TC:data:oscillations}, page~\pageref{subsec:TC:data:oscillations}) and present the results of this experiment compared to the results obtained when the oscillations are not removed. In this case, we find that removing the oscillations actually damages the performance of the EWS techniques, probably because by subtracting a periodic function from the entire time series we are adding a periodic function into the part of the time series close to the tipping point, where the periodicity in sea level pressure is temporarily lost. This is, then, a good test of the particular strength of the PS indicator, although the data available is, unfortunately, not sufficiently good quality to gain convincing results. In chapter~\ref{chap:2D} we include a separate discussion of the periodic component in the near-homoclinic system (System A) which orbits around a stable point before diverging when it crosses the saddle node (see section~\ref{sec:2D:periodic_orbit}, page~\pageref{sec:2D:periodic_orbit}). In this case we attempted two different methods to obtain a time series of one variable without the periodic component, but neither resulted in a clear EWS. We have concluded that the EWS indicator techniques are not well-suited to this particular system and so, in this case, the Jacobian-eigenvalue-estimation technique (see section~\ref{sec:2D:Williamson}, page~\pageref{sec:2D:Williamson}) is more appropriate, as we see in the results of this study (see section~\ref{subsec:2D:MethodsApplied:Williamson_and_ACF1DFA}, page~\pageref{subsec:2D:MethodsApplied:Williamson_and_ACF1DFA}). 
	\item \label{rep:gen:mathematical} \emph{The mathematical reasons behind the expected scaling behaviour should be more clearly expressed and discussed. }\\
	— We agree that this discussion was severely lacking in the original submission and most of chapter~\ref{chap:1D} is now dedicated to examining the power spectrum scaling exponent in more detail, particularly the mathematical reasons behind the expected scaling behaviour, and also the reasons for the choice of the range of frequencies in which the PS exponent is estimated. We begin by including the Wiener-Khinchin theorem in the section where the PS exponent is introduced (see page~\pageref{thm:1D:wiener-khinchin}), then we have expanded the sections where the relationships between the different scaling exponents are examined, in particular the relationship between the PS exponent $\beta$ and the ACF scaling exponent $\gamma$ (section~\ref{subsec:1D:relationships:beta_and_gamma}, page~\pageref{subsec:1D:relationships:beta_and_gamma}). Then, following the explanation of the \emph{critical slowing down} phenomenon, which is modelled by an AR(1) process $z_n = \mu z_{n-1} + \eta_n$ with increasing parameter $\mu$, we go on to show that by estimating the PS exponent $\beta$ of an AR(1) process in the frequency range $\m2\leq\log f\leq\m1$ the value is sensitive to change in $\mu$ due to the relation
	\begin{equation*}
	\beta = \log\left[  \frac{1+\mu^2-2\mu \cos(0.2\pi)}{1+\mu^2-2\mu \cos(0.02\pi)}  \right] \,,
	\end{equation*}
	(see equation~\ref{eqn:1D:PSindicator:PSofAR1}, page~\pageref{eqn:1D:PSindicator:PSofAR1}) and is therefore well-suited –in theory– to detecting critical slowing down, although we note that the power spectrum of the AR(1) process does not exhibit true asymptotic power-law scaling. The usefulness of estimating the PS exponent in the absence of power-law scaling is then examined further in section~\ref{subsec:1D:Exponents_as_EWS:nopowerlawsacaling} (page~\pageref{subsec:1D:Exponents_as_EWS:nopowerlawsacaling}) and these ideas are then used to establish the ``$\m2\leq\log f\leq\m1$'' frequency range in section~\ref{subsec:1D:Exponents_as_EWS:determining_frequency_range} (page~\pageref{subsec:1D:Exponents_as_EWS:determining_frequency_range}). The following sections (sections~\ref{subsec:1D:Exponents_as_EWS:PSsensitivity}, \ref{subsec:1D:Exponents_as_EWS:PSnonstationary} and \ref{subsec:1D:Exponents_as_EWS:Periodicity}, pages~\pageref{subsec:1D:Exponents_as_EWS:PSsensitivity}-\pageref{sec:1D:numerical_integration}) go on to investigate the sensitivity of the PS indicator to periodicity, changes in the AR(1) parameter and window size, particularly noting the effect of each on the scaling behaviour of the AR(1) process power spectrum. We believe that the reasons behind the use of the PS indicator as a measure of scaling behaviour (namely, that it is sensitive to critical slowing down), have now been established prior to the use of the PS indicator in toy models in section~\ref{sec:1D:indicatorComparison}.
	\item \label{rep:gen:periodogram} \emph{A definition of the periodogram is needed.}\\
	— The concept of the periodogram as an estimate of the power spectrum is now introduced briefly in the introductory chapter (see equation~\ref{eqn:intro:power_spectrum_def}, page~\pageref{eqn:intro:power_spectrum_def}), but a proper definition is also now given in chapter~\ref{chap:1D} (see definition~\ref{def:1D:PowerSpectrum}, page~\pageref{def:1D:PowerSpectrum}) where we additionally define the power spectrum scaling exponent and the method we use for its estimation. The Matlab function used to generate the fast Fourier transform periodogram and estimate the PS scaling exponent is included in appendix~\ref{appendix:PSexponent} (page \pageref{appendix:PSexponent}).
	\item \label{rep:gen:stochastic} \emph{The thesis is unclear regarding the way noise is implemented in the dynamical systems of interest. ... The author's understanding of the numerics of stochastic equations and Brownian motion is weak and must be reinforced.}\\
	— In the original submission the noise was integrated slightly differently depending on whether the system was in one variable or more than one. This was due simply to the progression of the work: different Matlab scripts were written as needed. We have now updated the thesis so that all dynamical systems are integrated numerically using the Milstein method described in section~\ref{sec:1D:numerical_integration} (page~\pageref{sec:1D:numerical_integration}). When re-integrating the systems used as examples, and then re-calculating the EWS indicators, we found that there was little or no difference to the original implementations, but we now have a consistent approach used throughout the whole thesis. The Matlab function of the Milstein method that we use in all our numerical examples is also included in appendix~\ref{appendix:milstein} (page~\pageref{appendix:milstein}).
	\item \label{rep:gen:TC} \emph{The application presented in chapter~\ref{chap:cyclones} (tropical cyclone) does not correspond to the definition of a tipping point. That is, a tropical cyclone does not conform to a general bistable system.}\\
	— We agree that the tropical cyclone probably does not conform to a general bistable system. Indeed, the model of the system that we develop in section~\ref{sec:TC:Model} (page~\pageref{sec:TC:Model}) is of a stationary noise process (with added periodicity) which is subject, at some point, to an exponential trend. However, we characterise this as a tipping point nonetheless, since it conforms to the definition ``an abrupt qualitative change in a dynamical system''. In our model the tipping point is \emph{transitional}, rather than \emph{bifurcational}. The definitions and varieties of tipping points are now discussed further in chapter~\ref{chap:intro} (see sections~\ref{subsec:intro:tippingpoints:defining_a_tipping_point} \& \ref{sec:intro:classification_of_TPs}, pages~\pageref{subsec:intro:tippingpoints:defining_a_tipping_point} \& \pageref{sec:intro:classification_of_TPs}). Some of the reasons why the tropical cyclone example was chosen are explained at the start of chapter~\ref{chap:cyclones}, and reiterated in the conclusions chapter. Briefly, the system was chosen because it gives an example of a system with an obvious tipping point (in the general sense of a sudden change in the system) but one that is probably not bifurcational, which allows us to do something new with the EWS indicator methods. In addition, the discovery that the sea-level pressure contains short-term oscillations (due to tidal effects) gives a good opportunity to study the effectiveness of the PS indicator (which we believe is insensitive to these kinds of periodicities) as compared to the lag-1 autocorrelation (which is not useful in oscillating systems). However, sufficiently good-quality data was not available to properly study these effects. This was an unfortunate issue and we comment further on this in the conclusions chapter (chapter~\ref{chap:Conc}).
	\item \label{rep:gen:99vs100} \emph{Some of the results of chapter~\ref{chap:cyclones} (e.g. the sensitivity of the results to using $99$vs.$100$ points as a window) are unconvincing.}\\
	— We agree that the argument for the presence of EWS was based on ad hoc analysis, and although a small trend in the PS indicator may be observed in the ensemble mean, it is far below any reasonable confidence level. This is commented upon in the results of chapter~\ref{chap:cyclones}, and again in the conclusion (chapter~\ref{chap:Conc}). 
	It is unfortunate that a system with mode convincing results was not available (despite our attempts to obtain, for instance, ``greening of the Sahara'' land model data from colleagues at the Max Planck Institute for Meteorology —see \citep{claussen2003climate,bathiany2013p1}— which is still not in public domain at sufficiently high time-resolution). Other suggestions which may be considered in future work are given in the conclusion, as are the reasons why the tropical cyclone example was chosen, which are detailed in the reply to comment~\ref{rep:gen:TC}. To answer the specific point about ``$99$vs.$100$ points as a window'', we cannot comment specifically on what causes this, and draw no conclusions of our own. This is simply an effect that we noticed and draw to the readers attention. We see in figure~\ref{fig:TC:PS_sensitivity} (page~\pageref{fig:TC:PS_sensitivity}) that there is a definite qualitative difference in the behaviour of the PS indicator applied to this data when using $\geq100$ points as opposed to using $<100$ points. This same effect is not observed when using the DFA indicator. Then again, in figure~\ref{fig:TC:stdSensitivity} (page~\pageref{fig:TC:stdSensitivity}) we see that while the standard deviation of the DFA and ACF1 indicators decreases steadily as the window size is increased, the standard deviation of the PS indicator, which appears also to decrease smoothly everywhere else, makes a sudden jump at exactly $100$ points. The results in figure~\ref{fig:TC:stdSensitivity} are generated separately, using a different Matlab script, to those in figure~\ref{fig:TC:PS_sensitivity}, but the ``99vs.100'' is observed in both. We have reviewed our method for estimating the PS exponent (a Matlab function, see appendix~\ref{appendix:PSexponent}) but there is no explanation here for the observed effect, and we note that with all of the data from the toy models used throughout the study (where the same Matlab function is used) there is no unexpected difference when using a window of 99 points as opposed to 100 points. We have also contacted the team at the UK Met Office who are responsible for these datasets, suspecting that some truncation or filtering algorithm may be responsible, but although they were happy to engage with us over the problem, we found no solution. We have left our observation as a side-note in chapter~\ref{chap:cyclones}, and simply state that this is, unfortunately, an unresolved issue.
	\item \label{rep:gen:conclusion} \emph{The thesis is in need of a final chapter which will give a summary of the results, put them into a broader perspective and propose future research.}\\
	— A concluding chapter is now included (see page~\pageref{chap:Conc}) in which the thesis is briefly summarised, conclusions are drawn from each separate line of the work, and suggestions for future work are made. The discussion sections at the ends of chapters~\ref{chap:1D}, \ref{chap:2D} and~\ref{chap:cyclones} have been left in since they aid the reading of the thesis as a complete work and are slightly different in tone from the sections of the concluding chapter.  
\end{enumerate}

\section*{Reply to preliminary report 1}
\label{rep:1}

The preliminary report of Examiner 1 makes three general comments to which we have made a reply below. Besides these there are also 46 specific comments. Some of these relate to small typos which have been silently corrected (all changes from the original manuscript appear in {\Red{red}} in this draft), but most require a short answer and these are supplied after the replies to the three general comments below.

\subsection*{General comments}

\begin{enumerate}
	\item  {\bf Comment:} Two of the early warning indicators used in the thesis are based on "scaling exponents", the PS and DFA scaling exponents. In addition, the scaling exponent of the auto-correlation function is also considered. Shortly, scaling is linked to power-law behaviour in which there is no specific characteristic scale such that $f(x)=f(cx)/c^\alpha$ where $\alpha$ is the scaling exponent. Unfortunately, this seems to be overlooked, as the scaling exponent was estimated even when it is clear that there is no scaling, like when there is crossover (AR(1) process) and when the time series is "oscillatory" (like the Van-der-Pol oscillator). Indeed, it is possible to estimate an exponent in a range where there seems to be a power-law behaviour but, from what I understood, the estimation of the scaling exponent was done without checking if indeed this is the case. Even in the case of tropical cyclones where the author noted that there is a clear tidal 12 hours frequency and even proposed ways to filter out this frequency, eventually this frequency was not filtered out. In the last case, it is not clear that indeed there is scaling (probably not). The above may underlies the weak performance of the indicators.\\
	{\bf Reply:} In the original submitted manuscript there was little explanation of the mathematical reasoning for using the PS exponent as an EWS indicator. We have attempted to rectify this by including several new sections in chapter~\ref{chap:1D} where we analyse the power spectra of different processes, most relevant of which is the AR(1) process mentioned in this comment. In our work we assume a tipping point (most notably a bifurcation) to be preceded by \emph{critical slowing down}, which is modelled as an AR(1) process $z_{n} = \mu z_{n-1} + \eta$ with increasing parameter $\mu$, as we explain in section~\ref{subsec:1D:Exponents_as_EWS:CriticalSD} (page~\pageref{subsec:1D:Exponents_as_EWS:CriticalSD}). It is therefore this increase in $\mu$ that must be detected by any EWS technique, and the lag-1 autocorrelation is an obvious choice. We show that although the power spectrum of the AR(1) process contains a crossover and does not exhibit true asymptotic power-law scaling, the \emph{estimated PS exponent}, which is obtained by a linear fit to the periodogram, is sensitive to change in $\mu$. Indeed, we derive equation~\ref{eqn:1D:PSindicator:PSofAR1} (page~\pageref{eqn:1D:PSindicator:PSofAR1}):
	\begin{equation*}
	\beta = \log\left[  \frac{1+\mu^2-2\mu \cos(0.2\pi)}{1+\mu^2-2\mu \cos(0.02\pi)}  \right] \,,
	\end{equation*}
	which relates the estimated PS exponent $\beta$ to the AR(1) parameter $\mu$. This equation is verified by numerical experimentation, the results of which are presented in figure~\ref{fig:1D:PSofAR1} (page~\pageref{fig:1D:PSofAR1}). The general justification for estimating the PS exponent of a power spectrum that does not exhibit true power-law scaling is then explained further in section~\ref{subsec:1D:Exponents_as_EWS:nopowerlawsacaling} (page~\pageref{subsec:1D:Exponents_as_EWS:nopowerlawsacaling}) and more discussion of the PS indicator and the AR(1) process, including the estimation of the PS exponent in oscillating time series, is presented in the following sections (see pages~\pageref{subsec:1D:Exponents_as_EWS:determining_frequency_range} to~\pageref{sec:1D:numerical_integration}). \\
	In the case of the tropical cyclone example in chapter~\ref{chap:cyclones}, we now give a more detailed discussion of the method used to remove the periodicity (see section~\ref{subsec:TC:data:oscillations}, page~\pageref{subsec:TC:data:oscillations}) and the results obtained from this experiment are now included in the manuscript. 
	\item {\bf Comment:} The term "potential" is frequently used in the thesis. It seems to be used in an inaccurate way. In physics "potential" is tightly linked to force through the relation $F=-\bigtriangledown(U)$ where the force is proportional to the second derivative of the location vector. In this context, a particle can be trapped in a potential well if its energy is smaller than the barrier's energy. However, in the examples considered in the dissertation, the potential is proportional to the first (time) derivative of the potential variable, e.g., $dx/dt=-dU/dx$. The "well" of U is not a real well as there is no relevant energy below which the "particle" is trapped. I would suggest the author to use the term "pseudo-potential" and to clarify that the potential he uses is not the same as the one used in physics in which a "well" is a well and a "barrier" is a barrier.\\
	{\bf Reply:} All instances of the term ``potential function'' in the manuscript have now been changed to ``generalised potential function'', consistent with the phrasing in \cite{nicolis2007}, etc., because of the similarity to the concept of an energy potential. References to \emph{wells}, as in the `double well' system, have been retained simply as illustrative descriptions, although we have taken care to explain in the manuscript that this is not related to the physical energy ``potential''. The local maxima in such functions are referred to as \emph{peaks}, rather than \emph{barriers}. This is explained in the manuscript (see page~\pageref{pageref:intro:generalisedpotential}).
	\item {\bf Comment:} It is not clear that indeed tropical cyclones belong to the class of tipping point mechanism. Please explain how this major example related to the tipping point discussed earlier.\\
	{\bf Reply:} Please refer to the reply to comment~\ref{rep:gen:TC} of the reply to the general report. 
	
\end{enumerate}

\subsection*{46 specific comments}

\begin{itemize}
	\item \emph{Comment 3} relates to the definition of the Fourier transform. — We have now properly defined the Fourier transform as we use it in the calculation of the power spectrum. See definition~\ref{def:1D:PowerSpectrum}, page~\pageref{def:1D:PowerSpectrum}.
	\item \emph{Comment 4} refers to the power spectrum of the white noise series, which is a constant value as a function of frequency. — We have now included a more detailed discussion of this on page~\pageref{def:1D:whitenoise}. The calculation using the Wiener-Khinchin theorem is given in equation~\ref{eqn:1D:PSofWhiteNoise} (page~\pageref{eqn:1D:PSofWhiteNoise}).
	\item \emph{Comment 5} suggests that the relationship between the PS and ACF exponents should be included in the discussion besides a derivation of why the PS exponent of red noise is $\m2$. — We have now included the Wiener-Khinchin Theorem on page~\pageref{thm:1D:wiener-khinchin} and the derivation of the red noise power spectrum on page~\pageref{eqn:1D:rednoise_scalingfactor}.
	\item \emph{Comment 6 (and 27)} refers to the dynamical system defined by a 2D system of ODEs, but which is not truly coupled and could be expressed as a single ODE. — In these specific cases we are using a system presented in another paper and attempting to reproduce the same results.
	\item \emph{Comment 12} refers to logarithmic binning. — We now explain the logarithmic binning applied to the periodogram prior to the estimation of the PS exponent (see page~\pageref{pageref:1D:logbinning}).
	\item \emph{Comments 13 and 16} refer to the use of two different methods for generating a red noise series. — The discrepancy here has now been corrected. We note that the AR(63) model and the AR(1) model are actually equivalent when pure red noise is generated (see page~\pageref{eqn:1D:ColouredNoise_kasdin}).
	\item \emph{Comment 15} refers to the range of frequencies, $10^{\m2}\leq f \leq 10^{\m1}$, in which the PS exponent is estimated. — The reasons for this choice are now explained thoroughly in section~\ref{subsec:1D:Exponents_as_EWS:determining_frequency_range} (page~\pageref{subsec:1D:Exponents_as_EWS:determining_frequency_range}).
	\item \emph{Comment 16} Refers to table~\ref{tab:1D:exponent_relationships}. — The error has now been corrected.
	\item \emph{Comment 17} Refers to the relationship between $\beta$ and $\gamma$, the PS and ACF scaling exponents. — We have now presented this relationship in more detail (see equation~\ref{eqn:1D:beta-gamma_linear_relation}, page~\pageref{eqn:1D:beta-gamma_linear_relation}).
	\item \emph{Comments 18 and 19} suggests that the deviations in figure~\ref{fig:1D:ExponentRelationships} (page~\pageref{fig:1D:ExponentRelationships}) can be explained by the fact that there is not true scaling for an AR(1) model and that there is crossover in the power spectrum. — This problem has been thoroughly explained in section~\ref{subsec:1D:Exponents_as_EWS:nopowerlawsacaling} (page~\pageref{subsec:1D:Exponents_as_EWS:nopowerlawsacaling}) and the discussion is expanded in the following sections (see pages~\pageref{subsec:1D:Exponents_as_EWS:determining_frequency_range} to \pageref{sec:1D:numerical_integration})
	\item \emph{Comments 21 and 23} call attention to the fact that the relationships between the PS, ACF and DFA exponents are only valid when there is scaling. — This is of course true, and we now comment upon this (see section~\ref{sec:1D:relationships}, page~\pageref{sec:1D:relationships}).
	\item \emph{Comment 22} suggests re-plotting figure~\ref{fig:1D:ArtificialSignal} with a continuously varying parameter, to avoid `jumps'. — As is the case for \emph{comment 6}, this system is chosen for the purpose of reproducing published results, and so the same system is used as accurately as possible. 
	\item \emph{Comment 24} makes the point that the constant term in the AR(1) model introduces a trend that will affect the calculation of the indicators. — This is very likely. We now provide a discussion of the effect of trends upon the ACF1 (and PS) indicator in section~\ref{subsec:1D:Exponents_as_EWS:PSsensitivity} (page~\pageref{subsec:1D:Exponents_as_EWS:PSsensitivity})
	\item \emph{Comment 26} makes the point that using the whole time series for the EOF calculation may result in bias since the length increase as the series becomes longer. — This is true, and has been commented upon in the text. In chapter~\ref{chap:cyclones} when the EOF method is applied to real data, we use the sliding-window EOF.
	\item \emph{Comment 28} points out that the DFA exponent has no meaning when applied to an oscillating series. — An investigation of the effect of periodicity on the DFA is now provided in section~\ref{subsec:1D:Exponents_as_EWS:Periodicity} (page~\pageref{subsec:1D:Exponents_as_EWS:Periodicity}). In this case the DFA indicator is simply presented as a comparison and we do not draw any conclusions from this result. 
	\item \emph{Comment 29} asks why we have used a diagonal transform matrix and put the coupling in the noise terms (non-diagonal noise matrix). — This is simply to aid the proceeding calculations. A system with non-diagonal transform matrix can be made into this form by a change of basis which will not affect the `blow up' behaviour. A calculation for the non-diagonal case is given in appendix~\ref{appendix:EOF-justification} (page~\pageref{appendix:EOF-justification}).
	\item \emph{Comment 30} comments on the high value of the DFA exponent in figure~\ref{fig:2D:AltEOFDFA_homoclinic} (page~\pageref{fig:2D:AltEOFDFA_homoclinic}), the same as figure~\ref{fig:2D:HomoclinicIndicators} (page~\pageref{fig:2D:HomoclinicIndicators}). — These experiments have been repeated in both cases and the same value is obtained. 
	\item \emph{Comments 31 and 37} ask if the Kendall rank coefficient in equation~\ref{eqn:2D:KendallCoef} should be normalised. — It should, and this has been updated.
	\item \emph{Comment 32} questions the reference to a tropical cyclone as a `tipping point'. — We consider that, although not bifurcational, the event is a \emph{transitional} tipping point. We now provide a discussion of the different types of tipping point in section~\ref{sec:intro:classification_of_TPs} (page~\pageref{sec:intro:classification_of_TPs}) and explain our choice of the tropical cyclone event in chapter~\ref{chap:cyclones} (see page~\ref{chap:cyclones}) and in the concluding chapter.
	\item \emph{Comment 33} requests the inclusion of the periodogram on a log-log plot so that evidence of scaling may be noted. — This is now included in section~\ref{subsec:TC:evidenceOfScaling} (page~\ref{subsec:TC:evidenceOfScaling}), in particular see figure~\ref{fig:TC:PS_loglog_RO12} (page~\pageref{fig:TC:PS_loglog_RO12}). We note that there is little evidence for true power-law scaling but that a general negative gradient in the power spectra can be observed.
	\item \emph{Comment 34} notes that the error in the PS exponent is very large in this application. — We agree that this is the case and consider this when reporting our results and drawing conclusions.
	\item \emph{Comment 35} remarks that very short series lengths cannot be used to obtain the DFA or PS exponent (referring to the range of window sizes as low as 40 points used in the sensitivity analysis). — The sensitivity analysis concludes that using 100 points or more is preferable, but we also find it interesting to note that the same pattern in the DFA indicator persists down to around 75 points when only five or six segment sizes are used for the DFA estimation.
	\item \emph{Comment 38} questions the presence of spikes in the value of the PS indicator at certain times when using certain EOF weighting schemes (figure~\ref{fig:TC:weightSensitivity}, page~\pageref{fig:TC:weightSensitivity}). — We agree that this is problematic and we cannot identify a specific cause. We only conclude that the EOF method is not well suited to this specific problem. 
	\item \emph{Comment 39} asks why one would choose to use EWS indicator methods for tropical cyclones, since there are good data-driven models that predict cyclones many days in advance. — We now explain in chapter~\ref{chap:cyclones} and in the conclusion (chapter~\ref{chap:Conc}) that our reasons for choosing this problem were motivated by the opportunity to test these methods in a novel scenario. First, the sea-level pressure contains oscillations which we hypothesise will negatively affect the ACF1 and DFA indicators whilst not affecting the PS indicator. Second, the tipping is non-bifurcational, which may present an interesting challenge. We refrain from making claims that these EWS indicator methods might actually be used in forecasting or prediction.
	\item \emph{Comment 40} notes that there is little similarity between the nine hurricanes in figure~\ref{fig:TC:3by3_sensitivity} and therefore it is unclear whether there will be a common EWS. — We agree with this comment and reply that we are looking only for some general trait, to which there may be exceptions. We see in the following figure (figure~\ref{fig_ContourMapSummary}, page~\pageref{fig_ContourMapSummary}) that the cyclones are generally moving from areas with a positive PS indicator gradient to areas with a lower or negative gradient. These results, however, are unconvincing, as we note in the conclusion (chapter~\ref{chap:Conc}). 
	\item \emph{Comment 41} asks about the origin of equation~\ref{eqn:TC:model:V_m}. — This equation is obtained by approximating the pressure profile of a hurricane as a rectangular hyperbola and using the `gradient wind equations' to obtain a formula for wind speed based on pressure, then discounting the Coriolis parameter as negligible in the region of maximum wind speed. We cite \cite{holland1980} which includes more details of the model equations. 
	\item \emph{Comment 42} asks if a mistake was made in equation~\ref{eqn:TC:ambientpressure}. — This has now been corrected. $K$ is in the range $[0,2\pi]$, not $[0,12]$.
	\item \emph{Comment 43} suggests that the results [of running the model] are expected since the statistics were tailored to that specific outcome. — This, in fact, was the purpose of constructing the model. We wish to included `realistic' noise into the deterministic model (realistic in that the PS exponent increases) without affecting the larger dynamics or the autocorrelation. 
	\item \emph{Comment 44} suggests the claim that the PS indicator ``clearly does begin to increase...'' is an exaggeration, since the increase is not significant. — This has now been rephrased (see page~\pageref{pageref:TC:exaggeration}).
	\item \emph{Comment 45} points out that the claim ``the PS and DFA exponents measure essentially the same thing'' is only valid for true scaling, which is not the case here. — This has now been rephrased (see page~\pageref{pageref:TC:scaling_comment}).
	\item \emph{Comment 46} refers to the lack of a separate concluding chapter. — See point~\ref{rep:gen:conclusion} of the reply to the general report. A concluding chapter is now included (see chapter~\ref{chap:Conc}, page~\pageref{chap:Conc}). 
\end{itemize}



\section*{Reply to preliminary report 2}
\label{rep:2}

Report 2 makes a number of specific suggestions. All except for points \ref{rep:2:toymodels} and \ref{rep:2:flag} are repeated in the general report (see page~\pageref{rep:gen}) and in these cases we refer back to the answers given to those comments. In the cases of points \ref{rep:2:toymodels} and \ref{rep:2:flag} we give a reply here. 

\begin{enumerate}
\item \emph{Joshua does not provide a definition of the periodogram in his thesis.} \\
— Please refer to point \ref{rep:gen:periodogram} of the reply to the general report. The definition of the periodogram is now included on page page~\pageref{def:1D:periodogram} of the thesis.
\item \emph{There is no mention of how integration of stochastic differential equations are performed.}\\
— Please refer to point \ref{rep:gen:stochastic} of the reply to the general report. Section \ref{sec:1D:numerical_integration} of the thesis (see page \pageref{sec:1D:numerical_integration}) now deals with this subject.
\item \label{rep:2:toymodels} \emph{Little discussion of the choice of parameters in many toy models (rate of change of parameters; intensity of noise).}\\
— We have now included a discussion of the choice of model parameters just before the presentation of those models and the results of the EWS analysis (see section~\ref{subsec:1D:indicatorComparison:choice_of_parameters}, page~\pageref{subsec:1D:indicatorComparison:choice_of_parameters}). Briefly, the parameters have generally been chosen so that the tipping point can be observed on timescales that will allow the EWS techniques to be applied, and so that the tipping dynamics are not obscured by the noise in the system.
\item \label{rep:2:flag} \emph{No clearly stated criteria for defining that the early warning indicator indeed flag the occurrence of a tipping point.} \\
— Unfortunately there is no EWS indicator technique where a consistently accurate `flag' of a tipping point exists. We have changed the wording of the manuscript in several places where the various indicator techniques are mentioned to make it clearer that we are not, indeed, attempting to identify such a flag. Such a thing would be specific to each different dynamical system, and would also be a function of the time-step used in the integration of the system, or the sampling rate of the measurements. In this thesis we are concerned simply with identifying early warning signals, usually an increasing value of the EWS indicator, but we do not study \emph{how} the indicator is increasing, nor are we able to say that when the indicator reached a given value, or a given rate of increase, that a tipping point will occur in $x$ time steps. We comment in the concluding chapter (chapter~\ref{chap:Conc}) that this is a long-term aim of the study of EWS indicators and that aspects of this thesis may be starting-points for future work in this direction.
\item \emph{The previous part describes the early warning indicators in relation to specific [...] bifurcations of low-dimensional systems [...yet...] the application takes place on data which have nothing to do with bifurcation, but which he still classifies as “tipping behaviour”, with no further discussion. I believe that a logical step is missing there.}\\
— Please refer to point \ref{rep:gen:TC} of the reply to the general report. A more detailed discussion of the definition of a `tipping point' is now included in the introductory chapter (see sections~\ref{subsec:intro:tippingpoints:defining_a_tipping_point} and \ref{sec:intro:classification_of_TPs}, pages~\pageref{subsec:intro:tippingpoints:defining_a_tipping_point} and \pageref{sec:intro:classification_of_TPs}).
\item \emph{Finally, the thesis does not have conclusions section.}\\
— Please refer to point \ref{rep:gen:conclusion} of the reply to the general report. A separate conclusions chapter is now included (chapter~\ref{chap:Conc}, page~\pageref{chap:Conc}).
\item \emph{...nor a sufficiently detailed discussion of tipping points in the Earth system.} \\
— Please refer to point \ref{rep:gen:debate} of the reply to the general report. A more detailed discussion of the literature is now included in the introductory chapter (see section~\ref{sec:intro:geophysical}, page~\pageref{sec:intro:geophysical}).
\end{enumerate}

\pagebreak